\chapter{The measure-theoretic vocabulary behind probability}
\label{ch:measure-vocabulary}
A probability model is not merely a formula for averages. It states
\emph{which outcomes exist}, \emph{which events can be assigned
probabilities}, and \emph{which calculations are allowed}.
This chapter provides the missing hypotheses for Results 21--24,
40--46, and 71--76.

\begin{definition}[Sigma-algebra and measure]
For a set $\Omega$, a $\sigma$-algebra $\mathcal F$ is a family of
subsets containing $\Omega$, closed under complements and countable
unions. A measure $\mu:\mathcal F\to[0,\infty]$ satisfies
$\mu(\varnothing)=0$ and countable additivity over disjoint events.
A probability measure has $\mu(\Omega)=1$. A measure is
\emph{$\sigma$-finite} if $\Omega$ is a countable union of
measurable sets of finite measure.
\end{definition}
\begin{definition}[Random variable and expectation]
A real-valued random variable $X$ is a measurable map
$(\Omega,\mathcal F)\to(\mathbb R,\mathcal B(\mathbb R))$:
for every Borel set $B$, its preimage $X^{-1}(B)$ is an event.
If $\int|X|\,d\Prob<\infty$, define $\E X=\int X\,d\Prob$.
The space $L^2$ consists of equivalence classes of variables with
$\E|X|^2<\infty$; equality almost surely is identified.
Its inner product is $\langle X,Y\rangle=\E[X\overline Y]$.
A \emph{Hilbert space} is a complete inner-product space; $L^2$
is one. Completeness is what makes orthogonal projection possible.
\end{definition}
\begin{proposition}[Conditional expectation is a projection]
For $X\in L^2$ and a sub-$\sigma$-algebra $\mathcal G\subseteq
\mathcal F$, the conditional expectation $\E[X\mid\mathcal G]$
is characterized by being $\mathcal G$-measurable and satisfying
$\E[(X-\E[X\mid\mathcal G])Z]=0$ for all bounded
$\mathcal G$-measurable $Z$. It is the orthogonal projection of
$X$ onto the closed subspace $L^2(\mathcal G)$.
\end{proposition}
\begin{proof}
The defining conditional-expectation identity gives orthogonality
for bounded $Z$. Truncation followed by $L^2$ approximation
extends it to $Z\in L^2(\mathcal G)$.
The projection theorem then characterizes the minimizer of
$\E[(X-Z)^2]$ among such $Z$.
\end{proof}
This explains why a conditional firing-rate prediction is an
optimal mean-square prediction \emph{relative to the variables
conditioned upon}, not relative to every unobserved brain state.

\begin{definition}[Stationarity, ergodicity, and predictability]
A process $X_t$ is \emph{wide-sense stationary} if its mean is
constant and $\Cov(X_t,X_s)$ depends only on $t-s$.
It is \emph{strictly stationary} if all finite-dimensional
joint laws are time-shift invariant. An invariant process is
\emph{ergodic} when invariant events have probability zero or one.
A predictable intensity $\lambda(t\mid\mathcal H_{t-})$ is
measurable relative to information available \emph{just before}
$t$; this forbids predicting a spike using its own occurrence.
The time-integrated intensity is called its \emph{compensator}
when the Doob--Meyer conditions are satisfied.
For a regular parametric model $p_\theta$, the \emph{score} is
$U_\theta=\nabla_\theta\log p_\theta(X)$ and
\emph{Fisher information} is
$\mathcal I(\theta)=\E_\theta[U_\theta U_\theta^\top]$,
under assumptions justifying its finiteness and calculus.
\end{definition}
\begin{proposition}[Law of total variance revisited]
For $X\in L^2$,
\[
\Var(X)=\E[\Var(X\mid\mathcal G)]
 +\Var(\E[X\mid\mathcal G]).
\]
\end{proposition}
\begin{proof}
Write $X-\E X=(X-M)+(M-\E X)$ where
$M=\E[X\mid\mathcal G]$. The cross term vanishes by
orthogonality. Taking squared expectations yields the identity.
\end{proof}
\begin{example}[LFP recording versus ensemble]
If a stationary and ergodic model describes a long LFP trace,
its time average can estimate the ensemble mean under integrability.
Stationarity alone does not justify equating those means, nor does
increasing the sample rate create independent trials.
\end{example}